% ---------------------------------------------------------------------
\chapter{Einrichten}
\label{kap:einrichten}

\begin{lernziele}
\begin{itemize}[nosep]
  \item Git so einrichten, dass jeder Commit korrekt dem eigenen Namen zugeordnet wird,
  \item die GitHub-Anmeldung auf der Kommandozeile einrichten,
  \item ein bestehendes Projekt klonen oder ein neues anlegen.
\end{itemize}
\end{lernziele}

\section{Die eigene Identität}

Jeder Commit speichert, \emph{wer} ihn erstellt hat. Darum muss Git einmalig Name und
E-Mail-Adresse kennen:

\begin{lstlisting}[style=befehle]
git config --global user.name  "Bahram Alizadah"
git config --global user.email "name@example.com"
\end{lstlisting}

\begin{teile}
  config        & Einstellungen von Git lesen oder schreiben \\
  --global      & gilt für alle Repositorys dieses Benutzers auf diesem Computer \\
  user.name     & der Name, der in jedem Commit als Autor erscheint \\
  user.email    & die E-Mail-Adresse des Autors \\
\end{teile}

Die E-Mail-Adresse sollte dieselbe sein wie im GitHub-Konto. Nur dann kann GitHub die
Commits dem Profil zuordnen und mit dem eigenen Bild anzeigen.

\begin{praxis}
In Unternehmen ist die korrekte Identität keine Formsache. Über die Geschichte lässt sich
für jede Codezeile nachvollziehen, wer sie wann eingeführt hat. Das hilft bei der
Fehlersuche („Wen frage ich, was hier gemeint war?“) und ist in regulierten Branchen wie
Banken oder Medizintechnik oft vorgeschrieben.
\end{praxis}

\section{Zwei sinnvolle Grundeinstellungen}

\begin{lstlisting}[style=befehle]
git config --global init.defaultBranch main
git config --global core.autocrlf true
\end{lstlisting}

\begin{teile}
  init.defaultBranch main & neue Repositorys erhalten den Hauptbranch \cmd{main}, wie es heute üblich ist \\
  core.autocrlf true      & passt unter Windows die Zeilenenden automatisch an (siehe Abschnitt~\ref{sec:zeilenenden}) \\
\end{teile}

Alle aktuellen Einstellungen zeigt:

\begin{lstlisting}[style=befehle]
git config --global --list
\end{lstlisting}

\section{Anmeldung bei GitHub}

Für Pull Requests auf der Kommandozeile gibt es das Programm \cmd{gh} (GitHub CLI). Es
muss einmal pro Computer angemeldet werden:

\begin{lstlisting}[style=befehle]
gh auth login --hostname github.com --git-protocol https --web
\end{lstlisting}

\begin{teile}
  auth login              & Anmeldung starten \\
  --hostname github.com   & bei welchem Dienst (es gibt auch firmeneigene GitHub-Server) \\
  --git-protocol https    & wie Daten übertragen werden; HTTPS ist der einfachste Weg \\
  --web                   & Anmeldung im Browser mit einem einmaligen Code \\
\end{teile}

Danach zeigt \cmd{gh} einen achtstelligen Code an. Im Browser wird er auf
\cmd{github.com/login/device} eingegeben und mit \emph{Authorize} bestätigt. \cmd{gh}
erhält dadurch einen \begriff{Token}, einen digitalen Schlüssel, mit dem es im Namen des
Benutzers handeln darf. Windows speichert ihn sicher in der Anmeldeinformationsverwaltung.

\begin{fehler}
In \emph{Git Bash} bricht \cmd{gh auth login} ohne Optionen mit „could not prompt“ ab.
Git Bash kann die Auswahlmenüs von \cmd{gh} nicht darstellen. Mit den oben gezeigten
Optionen entfallen die Menüs. Alternativ funktioniert PowerShell.
\end{fehler}

\section{Ein Repository beschaffen}

Es gibt zwei Wege zu einem Repository.

\textbf{Klonen} kopiert ein bestehendes Repository von GitHub auf den eigenen Computer,
samt vollständiger Geschichte. Das ist der Normalfall:

\begin{lstlisting}[style=befehle]
cd D:\github
gh repo clone gladtobethere/hslu-bachelor-ai-ml
\end{lstlisting}

Es entsteht der Ordner \texttt{D:\textbackslash{}github\textbackslash{}hslu-bachelor-ai-ml}, und \cmd{origin} zeigt
automatisch auf GitHub. Mit reinem Git lautet der Befehl \cmd{git clone <URL>}.

\textbf{Neu anlegen} macht einen bestehenden Ordner zu einem Repository:

\begin{lstlisting}[style=befehle]
git init                       # legt den Ordner .git an
\end{lstlisting}

\begin{erfahrung}
Ein Ordner pro Repository, an einem festen Ort, zum Beispiel \texttt{D:\textbackslash{}github\textbackslash{}<repo>}.
Mehrere Kopien desselben Repositorys auf einem Computer führen früher oder später dazu,
dass man in der falschen Kopie arbeitet und Änderungen „verschwinden“.
\end{erfahrung}

\begin{zusammenfassung}
\begin{itemize}[nosep]
  \item Einmalig \cmd{user.name} und \cmd{user.email} setzen, passend zum GitHub-Konto.
  \item \cmd{gh auth login} meldet die Kommandozeile bei GitHub an.
  \item Klonen holt ein bestehendes Repository, \cmd{git init} legt ein neues an.
\end{itemize}
\end{zusammenfassung}
